\section{Introduction}

Educational assistants must turn knowledge into instructional decisions: which
hint to offer after an unsuccessful attempt, what evidence would demonstrate
understanding, and how to organize practice within a limited timetable.
Producing a plausible explanation of a teaching method is easier than adapting
that method to a learner and delivering materials a teacher can use. This
distinction is central to pedagogical content knowledge and scaffolding
\citep{shulman1986,wood1976}, and remains relevant to language models:
MathTutorBench finds that subject-solving ability does not automatically
translate into effective tutoring \citep{mathtutorbench}.

Recent work addresses this gap through pedagogical instruction following and
explicit expert decisions \citep{learnlm,bridge}. Reusable \emph{Agent Skills}
offer a complementary approach: package a procedure, its templates, and
supporting resources so that an agent can consult them when needed
\citep{agentskills}. For example, an adaptive-hint package can connect evidence
of a student's difficulty to graduated assistance and a check for independent
understanding. Its educational value depends on whether the agent selects
applicable operations and completes them correctly. Merely making the document
available does not establish that connection.

Existing benchmarks provide complementary evidence. EduBench covers diverse
educational requests, while MathTutorBench and MRBench assess open-ended
tutoring behavior \citep{edubench,mathtutorbench,mrbench}. SkillsBench evaluates
procedural packages through paired agent runs and executable verification
\citep{skillsbench}. Open-ended teaching products introduce a further challenge:
a generated lesson, diagnostic question, or practice schedule may be fluent
and procedurally detailed yet infeasible or substantively incorrect. Measuring
Skill use in this setting requires linking a teaching task to its procedural
resource, the agent's actions, and the quality of the delivered product.

We introduce \bench{}, a unified benchmark of 305 tasks across 14 educational
Skill families. Each task binds a teaching situation, a designated Skill, and
applicable rubric criteria. Paired \noskill{} and \withskill{} evaluations
compare delivered outcomes with and without procedural support; native execution records
show whether agents actually load it. Reference-conditioned judgments assess
the delivered text provided to the evaluator, while a 28-task expert-adjudicated sample checks
how those judgments behave across the 14 families. A complementary Qwen3-4B
study examines learning from Skill-informed artifacts through supervised
fine-tuning (SFT) and subsequent outcome-based reinforcement learning.

The analysis reveals a gap between access and useful delivery. Across four
native-access models, expected-package loading ranges from 17.70\% to 71.48\%,
but aggregate With--No changes remain between $-1.19$ and $+0.50$ score points.
Family breakdowns and paired cases locate more explicit teaching operations,
resource failures, and incomplete deliverables behind those averages.
Human-adjudicated scores further show that close reference-relative scores can
coexist with substantial differences in their underlying quality components.
These findings motivate evaluating procedural support at both the execution
and artifact levels.

Our contributions are threefold: (1) a 305-task educational benchmark linking
14 reusable Skills to explicit output requirements and rubrics; (2) paired
model evaluations with family-level outcomes, observable Skill access, and
trace-based explanations; and (3) complementary studies of artifact learning,
expert-adjudicated evaluation, and execution cost. Together, they make the
path from procedural guidance to educational work empirically inspectable.
